\subsection{02.12.2025}

\subsubsection{Ziel}
Das Ziel des Tages war es, die restlichen Bibliotheken zu installieren
und anschließend die ersten Lektionen der JOY-PI-Anleitung durchzuführen.

\subsubsection{Durchgeführte Arbeiten}
Für den IR-Sensor wurde die Bibliothek
\texttt{IR-Remote-Receiver-Python-Module} von „owainm713“
installiert. Zusätzlich wurden weitere benötigte Bibliotheken in der
virtuellen Python-Umgebung installiert.

Anschließend wurde der Quellcode der Lektion 1 (Verwenden des Buzzers für Warntöne) der
JOY-PI-Anleitung analysiert. Danach wurden die Lektionen 2 (Buzzer mit Taste steuern) 
und 3 (Wie ein Relais funktioniert und wie man es steuert)
durchgeführt, um zu überprüfen, ob die Installation der Bibliotheken
erfolgreich war und um die Funktionsweise der GPIO-Pins besser zu
verstehen.

Zusätzlich wurden die Lektionen 4 (Senden Sie ein Vibrationssignal) und 5 
(Geräusche mit dem Schlalsensor erkennen) betrachtet, jedoch nicht aktiv
durchgeführt. Ziel war hierbei hauptsächlich das Verständnis des
Programmaufbaus und der verwendeten Funktionen.


\subsubsection{Problem 1: IR-Romote-Receiver-Python-Module}
Die Installation der Bibliothek für den IR-Sensor war zunächst auch nach
mehreren Versuchen nicht erfolgreich. Währenddessen wurde unter anderem
eine Anmeldung bei GitHub gefordert.

\textbf{Lösung}
Es stellte sich heraus, dass beim Abtippen des Installationsbefehls ein
Tippfehler entstanden war. Nach der Korrektur des Befehls konnte die
Bibliothek erfolgreich installiert werden.

\subsubsection{Erkentnisse}
Es wurde festgestellt, dass die in der Anleitung angegebenen Befehle zur
systemweiten Installation, beispielsweise
\texttt{sudo python3 install <package>}, nicht innerhalb der virtuellen
Umgebung verwendet werden sollten.

Stattdessen müssen die Bibliotheken innerhalb der aktivierten virtuellen
Umgebung mit dem Befehl \texttt{pip install <package>} installiert werden.

Über den Aufbau von Python-Code wurden folgende Erkenntnisse gewonnen:
\begin{itemize}
    \item Zur Nutzung der GPIO-Pins des Raspberry Pi muss die Bibliothek
    \texttt{RPI.GPIO} importiert werden.
    \item Einer Variable kann die Nummer eines GPIO-Pins zugewiesen
    werden, um diesen im Programm anzusprechen.
   \item Es muss ein GPIO-Modus festgelegt werden, damit die Pins korrekt
    angesprochen werden können.
    \subitem Beispiel:
    \texttt{GPIO.setmode(GPIO.BCM)}
    \item Es muss definiert werden, ob ein Pin als Input oder Output
    verwendet werden soll.
    \subitem Beispiele:
    \texttt{GPIO.setup(pinnummer, GPIO.IN, pull\_up\_down=GPIO.PUD\_UP)}
    \texttt{GPIO.setup(pinnummer, GPIO.OUT)}
    \item Output-Pins können mit
    \texttt{GPIO.output(pinnummer, GPIO.HIGH)} oder
    \texttt{GPIO.output(pinnummer, GPIO.LOW)}
    aktiviert beziehungsweise deaktiviert werden
    \item Mit \texttt{GPIO.input(pinnummer)} können Werte von Input-Pins
    abgefragt werden.
    \item Für gewollte Endlosschleifen wird häufig ein Try-Except-Block
    verwendet, damit die Schleife mit einem KeyboardInterrupt
    (\texttt{Strg + C}) beendet werden kann.
    \item Beim Verlassen des Programms sollte mit
    \texttt{GPIO.cleanup()} eine Rücksetzung der GPIO-Pins erfolgen,
    damit diese nicht aktiv bleiben.
\end{itemize}

Über Relays wurden folgende Erkenntnisse gewonnen:

\begin{itemize}
    \item Bei der Nutzung von NC + CO wird der Stromkreis geschlossen,
    wenn \texttt{GPIO.LOW} anliegt.
    \item Bei der Nutzung von NO + CO wird der Stromkreis geschlossen,
    wenn \texttt{GPIO.HIGH} anliegt.
\end{itemize}

\subsubsection{Nächste Schritte}
Als nächster Schritt sollen weitere Lektionen der JOY-PI-Anleitung
angesehen und gegebenenfalls durchgeführt werden, um das Verständnis
der GPIO-Pins und ihrer Ansteuerung weiter zu vertiefen.